\chapter{Related Work}
\label{RelatedWork}

This chapter surveys the literature that is most relevant to the thesis. The review is organised into three thematic strands. The first strand covers interval-based event detection and the ISEQL query language, which form the temporal reasoning foundation of the system. The second strand covers vision-language models and their application to video surveillance, including the immediate predecessor work. The third strand covers environmental sound classification and the use of audio in surveillance contexts. The chapter concludes by identifying the research gap that this thesis addresses.

\section{Interval-Based Event Detection}

The temporal representation of events is a fundamental problem in artificial intelligence and database systems. \citet{allen1983} introduced an interval algebra that defines thirteen possible binary relations between time intervals, including relations such as BEFORE, MEETS, OVERLAPS, DURING, and STARTS. These relations provide a formal calculus for reasoning about temporal descriptions of events and have been widely adopted in temporal databases, planning systems, and event detection frameworks.

Building on Allen's interval algebra, \citet{helmer2016highlevel} proposed ISEQL, a query language based on relational algebra extended with interval operators for detecting high-level surveillance events from a video stream. The language defines operators such as Left Overlap Join (LOJ), During Join (DJ), Before (BEF), and Start Preceding (SP), each parameterised by threshold values that control the temporal tolerance. The authors evaluated ISEQL on the ITEA CANDELA dataset, a train station benchmark, and a university parking lot dataset, achieving perfect precision and recall on most event classes with queries completing in under forty milliseconds for a full video.

\citet{persia2017interactive} extended this work with an interactive framework for video surveillance event detection and modeling. The framework adopted a three-layer architecture consisting of an image processing library, an interval action detector, and a high-level event detector. Each layer was exposed through a graphical user interface that allowed users to inspect the system's internal state, create event definitions on the fly, and track processing through each stage. The system supported both offline and online detection modes and was demonstrated on the CANDELA, PETS 2006, and university parking lot datasets.

In parallel, \citet{persia2017labeling} developed a framework for labeling video frame sequences with interval events at the medium level of abstraction. The system used OpenCV-based image processing to identify and track people and items, then exploited interval-based relational algebra for offline and online labeling. The medium-level predicates included IN (person entering a region), HasPkg (person carrying a package), and InsideCar (person entering or exiting a vehicle). The experimental evaluation achieved precision and recall of 1.0 on the IN predicate and high recall on the HasPkg predicate after threshold tuning.

\citet{piatov2021cache} developed a family of efficient plane-sweeping interval join algorithms for evaluating a wide range of interval predicates, including Allen's relationships and parameterised relationships. Their approach used a framework of composable components that could be flexibly combined to support different interval relations. By employing a compact gapless hash map data structure, the algorithms utilised the CPU cache efficiently and achieved several times faster performance than state-of-the-art techniques.

\section{Vision-Language Models for Surveillance}

Recent advances in vision-language models have opened new possibilities for video surveillance. VLMs are multimodal models trained on large datasets of image-text pairs that can generate natural language descriptions of visual content, answer questions about images, and reason about spatial relationships between objects. Foundational VLMs such as CLIP \citep{radford2021clip}, Flamingo \citep{alayrac2022flamingo}, BLIP-2 \citep{li2023blip2}, and LLaVA \citep{liu2024llava} have demonstrated that aligning visual features with language model embeddings enables zero-shot transfer to a wide range of vision tasks. More recent video-focused VLMs such as Video-LLaMA \citep{zhang2023videollama} extend this capability to video understanding, and comprehensive surveys of video LLMs \citep{tang2025vidllm} provide a systematic overview of the field. The present thesis applies these VLM capabilities specifically to the surveillance domain.

\citet{crescitelli2025design} introduced a hybrid framework that integrated VLM-based perception with the ISEQL interval-based query language for temporal reasoning. The framework consisted of four modules: a VLM module that processed each video frame through a spatial grid prompt to extract relation triples, a data adaptation module that converted VLM output into a temporal database format, an interval action detector that aggregated frame-level observations into intervals, and a high-level event detector that composed interval events using ISEQL temporal relationships. A key innovation was the use of ISEQL spatial relationships within the interval action detector to correct VLM errors caused by occlusion and missed detections. The framework was evaluated on the CANDELA, PETS 2006, university parking lot, and VIRAT \citep{oh2011virat} datasets, demonstrating that the structured pipeline outperformed direct full-event VLM prompting.

\citet{crescitelli2026vismode} presented VIS MODE, a four-page demo paper describing the pipeline as a complete system. The VIS MODE system evaluated several VLMs including Qwen3-VL, Gemini 2.5 Pro, and Gemini 3 Pro, with Qwen3-VL producing the best results. The paper reported one hundred percent detection on high-quality footage with an average time error of approximately 0.5 seconds, and eighty-seven percent precision and recall on low-quality legacy datasets. However, three methodological weaknesses limit the reproducibility of these claims. First, the system did not include object re-identification: each frame was processed independently, so the VLM assigned new identifiers to the same person or vehicle on every frame, and intervals never spanned multiple frames. Second, the evaluation was conducted on unlabeled ground truth, so the reported precision and recall could not be independently verified. Third, the system was non-deterministic: inference used a temperature greater than zero with no fixed seed, so the VLM freely generated descriptions that varied across runs. Because every call produced fresh identifiers, even the carried object received a new ID on each frame, so the handoff query could never match a single object across two persons and failed completely. In the authors' GATCHA implementation this combination of limitations resulted in zero detected events, since the tracking-dependent intervals never formed. The present thesis addresses all three weaknesses by adding a deterministic two-phase re-identification mechanism, evaluating on per-frame labeled ground truth, and fixing temperature at zero with seed 42 to ensure reproducible inference. The present thesis evaluates a broader set of VLMs including Pixtral 12B, Ministral 3-14B, Gemini 2.5 Flash, and Gemini 3.6 Flash.

\section{Environmental Sound Classification}

Environmental sound classification is the task of identifying non-speech sounds from audio recordings. Audio representation learning has progressed from supervised CNNs to self-supervised approaches such as wav2vec 2.0 \citep{baevski2020wav2vec2} and HuBERT \citep{hsu2021hubert}, and to Transformer-based architectures such as AST \citep{gong2021ast} and BEATs \citep{chen2023beats}. The CLAP model \citep{elizalde2023clap} learns a shared embedding space between audio and natural language, enabling zero-shot audio classification. More recently, large audio-language models (LALMs) have extended large language models to process audio input directly, with Qwen-Audio \citep{chu2023qwenaudio}, Qwen2-Audio \citep{chu2024qwen2audio}, and SALMONN \citep{tang2024salmonn} demonstrating general-purpose audio understanding capabilities. The LALM survey \citep{luo2026lalmsurvey} provides a comprehensive taxonomy of these models. \citet{kong2020panns} proposed PANNs (Pretrained Audio Neural Networks), a family of convolutional neural networks trained on the full AudioSet dataset of over five thousand hours of audio with 527 sound classes. The best-performing architecture, Wavegram-Logmel-CNN, achieved a state-of-the-art mean average precision of 0.439 on the AudioSet tagging task, outperforming the previous best system by 0.047. The authors demonstrated that PANNs could be transferred to six downstream audio pattern recognition tasks including ESC-50, speech emotion classification, and music classification, achieving state-of-the-art performance in several of them.

AudioSet itself was introduced by \citet{gemmeke2017audioset} as a large-scale dataset of manually annotated audio events organised in a hierarchical ontology covering human and animal sounds, musical instruments, and environmental noises. The ESC-50 dataset \citep{piczak2015esc50} provides a smaller but curated benchmark of 2000 environmental sound clips across 50 classes and is commonly used as a transfer learning evaluation target for audio models. AudioSet's ontology is well suited as a source for mapping to surveillance-relevant sound classes.

In the context of video surveillance, environmental sound offers a complementary signal to visual perception. While a VLM can only detect events within the camera's field of view, an audio classifier can detect sounds originating outside the frame, such as shouts, breaking glass, and vehicle noises from adjacent areas. Prior surveillance systems have primarily relied on speech-based audio processing using automatic speech recognition (ASR) and speech emotion recognition (SER), but these approaches depend on the presence of speech, which is not guaranteed in surveillance audio. The present thesis departs from this speech-centric paradigm by using PANNs CNN14 and Qwen2-Audio-7B-Instruct for environmental sound detection, which remain effective even in the absence of speech.

\section{Research Gap}

The existing literature establishes strong foundations in three areas: interval-based event detection through ISEQL, VLM-based visual perception for surveillance, and environmental sound classification through pretrained audio neural networks. Recent surveys on multimodal anomaly detection \citep{barbosa2025multimodalvad} highlight the growing interest in fusing multiple sensor modalities for surveillance, but they do not propose a specific interval-based temporal fusion mechanism. The intersection of interval temporal logic and LLMs has also begun to receive attention \citep{bellodi2025intervalitl}, yet no prior work has combined these three strands into a unified multimodal surveillance framework and systematically evaluated the contribution of each modality through controlled ablation.

The VIS MODE system \citep{crescitelli2026vismode} demonstrated that VLM-based perception combined with ISEQL temporal reasoning can detect complex surveillance events. However, it used visual input only. The present thesis extends VIS MODE by adding an audio stream and defining a three-condition ablation that isolates the visual, audio, and fused conditions. The ISEQL framework \citep{helmer2016highlevel, persia2017interactive} provided the interval operators and the architectural pattern, but did not incorporate deep learning based perception modules. The present thesis connects these interval operators to modern visual and audio backends.

Furthermore, prior work on audio in surveillance has focused on speech and emotion recognition, which requires the presence of speech. The present thesis uses environmental sound classification, which captures a wider range of surveillance-relevant events and does not depend on speech. The key research gap is the absence of a controlled ablation study that quantifies the marginal contribution of each modality on scenes where only one modality has signal, and demonstrates that multimodal fusion guards against the false positive rate of the audio-only condition.

This thesis addresses that gap by defining three conditions (visual-only, audio-only, multimodal UNION) evaluated on binary verdict accuracy across 30 labeled scenes. The ablation study answers the question of whether combining both modalities via UNION improves detection performance over either modality alone.
